\section{Introduction}
\label{sec:intro}

Inverse folding is solved well enough to be a commodity: ProteinMPNN~\cite{proteinmpnn}
has 1.66 M parameters (measured from the released weights; the paper reports
1.678 M in its own table), fits on a single A100, and recovers about 52\% of
native residues on its own PDB test set. The bottleneck has moved downstream ---
to whether the designed sequence can actually be made. Solubility, aggregation,
expression and electrostatic behaviour are now the binding constraints in real
campaigns, and the community's answer has been to \textbf{retrain the prior} on
filtered data: SolubleMPNN~\cite{soluble_mpnn} excludes transmembrane entries,
HyperMPNN~\cite{hypermpnn} trains on thermophile structures,
MembraneMPNN~\cite{membranempnn} conditions on a membrane label.

That answer has five consequences for evaluation practice, and this paper
measures each of them. The first three are properties of the models; the last
two are properties of how we evaluate them.

\paragraph{P1 --- single points.}
A retrained model returns one point in property space. There is no knob for
``20\% less hydrophobic than last time'', only a different checkpoint.

\paragraph{P2 --- unreported leakage.}
Moving one property moves others. SolubleMPNN, asked only to remove
transmembrane proteins, also shifts net charge by $-3.5$ units and pI by
$-1.2$. Neither its main text nor the follow-up literature we could find reports
this.

\paragraph{P3 --- blind acceptance metrics.}
Recovery, scTM and pLDDT do not see P2. In our data the six published models,
which disagree by 6.0 charge units, are indistinguishable on every standard
metric (scTM spans 0.78 to 0.88, and the spread is driven by one retrained model
rather than by the property edits).

\paragraph{P4 --- saturation.}
More than a third of the benchmark backbones are at ceiling for every model, so
those cells can only confirm agreement; and because arm differences are bimodal,
the arm mean is dominated by a handful of catastrophic cells. An evaluation that
reports only means therefore cannot tell ``uniformly slightly worse'' from
``fine everywhere except three catastrophic failures''.

\paragraph{P5 --- oracle dependence.}
The comparison is reported against one refolding oracle. We show that a second
oracle shifts absolute scTM by 0.15 and does not preserve the arm ranking, which
makes single-leaderboard reporting an artefact of the choice of oracle rather
than a property of the models.

Our contributions:

\begin{enumerate}
\item \textbf{An evaluation instrument} for control-style design methods: the
  controllability slope $k$, leakage, and the per-backbone tail, each with a
  protocol (backbone as the unit of analysis, paired bootstrap, Holm correction,
  and an explicit minimum detectable effect) that makes negative results
  interpretable.
\item \textbf{A controlled comparison of the whole published
  developability-MPNN family} (ProteinMPNN, SolubleMPNN, HyperMPNN,
  MembraneMPNN, LigandMPNN) against jointly-scored guidance arms on 35 backbones
  and 5,936 designs, with a second oracle.
\item \textbf{Five findings about current practice}, each tied to a measured
  failure: 4--13\% property movement, 6.0 units of inter-model disagreement
  invisible to standard metrics, bimodal tail behaviour hidden by means, a third
  of the benchmark saturated, and a ranking that does not survive a change of
  oracle. Two further results follow from measuring control on the full
  benchmark: the control-folding frontier is flat over a tenfold range of control
  strength, and the position-specific part of the reward is load-bearing (a
  global penalty controls the target better but costs 0.033 scTM and leaks
  charge).
\item \textbf{A closed-form guidance baseline}: for composition-linear
  properties the any-order-autoregressive guidance of
  ProteinGuide~\cite{proteinguide} reduces to a constant per-amino-acid logit
  bias, so a frozen ProteinMPNN becomes a calibrated, reversible knob at no
  extra predictor cost --- the cheapest possible control method, used here as
  the instrument that exposes what the published models' single points cannot
  show.
\end{enumerate}

\textbf{What is and is not new here.} The any-order-autoregressive guidance
identity is due to ProteinGuide~\cite{discrete_guidance,proteinguide} (Nisonoff
et al., 2025; Xiong et al., 2026); we do not claim it. Our contribution is (i)
the closed-form specialisation to composition-linear rewards, which needs no
property predictor at inference time, and (ii) the benchmark that measures what
such control costs, including the leakage and per-backbone analyses. Guidance
through a learned predictor for non-linear objectives (the TAG route in
ProteinGuide) is deliberately out of scope and is listed in the limitations; a
reader should not read this paper as a comparison of guidance algorithms.
